%------------------------------------------------------------------------------%
\documentclass[a4paper,12pt,fleqn]{article}

\usepackage{calculo_varias_variaveis-1}

\ShowAnswers

%------------------------------------------------------------------------------%
\begin{document}

\makeHeader{CFVVI}{Prova 3 -- Turma 2}{05/11/2025}

% 01
%------------------------------------------------------------------------------%
\exercise{50\%}
Determine a aproximação linear
de \(f(x, y) = e^{xy}\cos(xy)\) no ponto \((\pi, 1)\) e
use-a para aproximar \(f(3.1,\,1.02)\).
Utilize apenas 2 casas decimais.
\clearpagequestiononly

\begin{answer}
A aproximação linear de \(f\) em \((a, b)\) é dada por
\[
  L(x, y) = f(a, b) + f_x(a, b)(x-a) + f_y(a, b)(y-b)
\]
Cálculo das derivadas parciais
\begin{align*}
  f_x(x, y)
  & = \D{}{x}\left(e^{xy}\cos(xy)\right) \\
  & = \D{}{x}\left(e^{xy}\right) \cos(xy)
    + e^{xy} \D{}{x}\left(\cos(xy)\right) \\
  & = e^{xy}\D{}{x}\left(xy\right) \cos(xy)
      + e^{xy} \left(-\sin(xy)\right)\D{}{x}(xy) \\
  & = e^{xy}y \cos(xy) - e^{xy} \sin(xy)y \\
  & = ye^{xy}\left(\cos(xy)-\sin(xy)\right)
\end{align*}
\begin{align*}
  f_y(x, y)
  & = \D{}{y}\left(e^{xy}\cos(xy)\right) \\
  & = \D{}{y}\left(e^{xy}\right) \cos(xy)
    + e^{xy} \D{}{y}\left(\cos(xy)\right) \\
  & = e^{xy}\D{}{y}\left(xy\right) \cos(xy)
      + e^{xy} \left(-\sin(xy)\right)\D{}{y}(xy) \\
  & = e^{xy}x \cos(xy) - e^{xy} \sin(xy)x \\
  & = xe^{xy}\left(\cos(xy)-\sin(xy)\right)
\end{align*}
Avaliação em \((\pi, 1)\):
\[
    f(\pi, 1)
    = \left(e^{xy}\cos(xy)\right)\evalat{(\pi, 1)}{}
    = e^\pi\cos(\pi)
    = -e^\pi
\]
\begin{align*}
    f_x(\pi, 1)
    & = \left(ye^{xy}\left(\cos(xy)-\sin(xy)\right)\right)\evalat{(\pi, 1)}{} \\
    & = 1e^\pi\left(\cos(\pi)-\sin(\pi)\right)
    = e^\pi\left(-1-0\right)
    = -e^\pi
\end{align*}
\begin{align*}
    f_y(\pi, 1)
    & = \left(xe^{xy}\left(\cos(xy)-\sin(xy)\right)\right)\evalat{(\pi, 1)}{} \\
    & = \pi e^\pi\left(\cos(\pi)-\sin(\pi)\right)
    = \pi e^\pi\left(-1-0\right)
    = -\pi e^\pi
\end{align*}
A aproximação linear é
\begin{align*}
  L(x, y)
  & = f(a, b) + f_x(a, b)(x-a) + f_y(a, b)(y-b) \\
  & = f(\pi, 1) + f_x(\pi, 1)(x-\pi) + f_y(\pi, 1)(y-1) \\
  & = -e^\pi - e^\pi(x-\pi) - \pi e^\pi (y-1) \\
  & = -e^\pi \left[ 1 + (x-\pi) + \pi (y-1)\right]
\end{align*}
\begin{align*}
  L(3.1,\,1.02)
  & = -e^\pi \left[ 1 + (x-\pi) + \pi (y-1)\right] \evalat{(3.1,\,1.02)}{} \\
  & = -e^\pi \left[ 1 + (3.1-\pi) + \pi (1.02-1)\right] \\
  & = -e^\pi \left[ 1 + (-0.04) + \pi (0.02)\right] \\
  & = -e^\pi \left[ 0.96 + 0.02\pi \right]
\end{align*}
Com uma calculadora podemos avaliar que
\begin{align*}
  L(3.1,\,1.02) = -23.669 \\
  f(3.1,\,1.02) = -23.613
\end{align*}
\end{answer}

% 05
%------------------------------------------------------------------------------%
\exercise{50\%}
Considere a função \(f:\mathbb{R}^3\to\mathbb{R}\) dada por
\(
  f(x, y, z) = x e^{yz} + \sin(yz)
\)
Calcule a derivada direcional de \(f\) no ponto \(P = (1, 0, 2)\) na direção do vetor
\(\mathbf{v}=(2, 1, -1)\).

\begin{answer}
Normalização o vetor \(\mathbf{v}\)
\[
  \|\mathbf{v}\| = \sqrt{2^2+1^2+(-1)^2}=\sqrt{6},
  \qquad
  \mathbf{u} = \frac{1}{\sqrt{6}}(2, 1, -1).
\]
Gradiente de \(f\):
\begin{align*}
  f_x(x, y, z) & = e^{yz}                \\
  f_y(x, y, z) & = xz e^{yz} + z\cos(yz) \\
  f_z(x, y, z) & = xy e^{yz} + y\cos(yz)
\end{align*}
Avaliação no ponto \(P=(1,0,2)\)
\begin{align*}
f_x(1, 0, 2) & = e^{0} = 1 \\
f_y(1, 0, 2) & = 1 \cdot 2 \cdot e^{0} + 2\cos(0) = 2+2 = 4 \\
f_z(1, 0, 2) & = 1 \cdot 0 \cdot e^{0} + 0\cos(0) = 0
\end{align*}
Assim,
\(
  \nabla f(1, 0, 2) = (1, 4, 0)
\)

A derivada direcional é o produto escalar
\begin{align*}
  D_{\mathbf{u}}f(1, 0, 2)
  & = \nabla f(1, 0, 2)\cdot \mathbf{u} \\
  & = (1, 4, 0)\cdot \frac{1}{\sqrt{6}}(2, 1, -1) \\
  & = \frac{1}{\sqrt{6}}(1\cdot 2 + 4\cdot 1 + 0\cdot (-1)) \\
  & = \frac{6}{\sqrt{6}} \\
  & = \sqrt{6}
\end{align*}.
\end{answer}

%------------------------------------------------------------------------------%
\end{document}
%------------------------------------------------------------------------------%